\section*{Chapitre \ref{ch:b2:plant-plasticity} --- Adaptations des plantes et plasticité phénotypique}

\begin{solution}{exo:b2:plant-plasticity:1}
Plasticité~: un \omterm{def:b2:meiosis-heredity:vocabulary}{génotype}, plusieurs \omterm{def:b2:meiosis-heredity:vocabulary}{phénotypes} selon l'environnement
(\omterm{prop:b2:plant-plasticity:sunshade}{feuilles de lumière} et d'ombre d'un même chêne). \omterm{def:b2:plant-plasticity:plasticity}{Norme de réaction}~: la
courbe du \omterm{def:b2:meiosis-heredity:vocabulary}{phénotype} d'un \omterm{def:b2:meiosis-heredity:vocabulary}{génotype} en fonction d'une variable de
l'environnement (épaisseur des feuilles en fonction de la lumière).
\omterm{def:b2:plant-plasticity:plasticity}{Acclimatation}~: un ajustement physiologique réversible (endurcissement
au gel du seigle en automne). \omterm{def:b2:plant-plasticity:plasticity}{Adaptation}~: un ajustement héréditaire
sélectionné au fil des générations (le métabolisme CAM d'un cactus).
\end{solution}

\begin{solution}{exo:b2:plant-plasticity:2}
Coléoptile intact éclairé d'un côté~: il se courbe vers la lumière.
Extrémité couverte d'un capuchon opaque~: pas de courbure --- c'est
l'extrémité qui perçoit la lumière. Extrémité couverte d'un capuchon
transparent~: courbure --- le capuchon lui-même est sans effet. Base
protégée par un tube opaque, extrémité exposée~: courbure --- la
perception a lieu à l'extrémité et la réponse en dessous.
\end{solution}

\begin{solution}{exo:b2:plant-plasticity:3}
\omterm{prop:b2:plant-plasticity:sunshade}{Feuille de lumière}~: petite, épaisse, deux ou trois assises
palissadiques, cuticule épaisse, nombreux stomates, photosynthèse
maximale élevée et point de compensation élevé. \omterm{prop:b2:plant-plasticity:sunshade}{Feuille d'ombre}~:
grande, fine, une seule assise palissadique, plus de chlorophylle par
gramme, faible respiration, point de compensation bas. Le choix est fait
lorsque la feuille est un primordium dans le bourgeon, d'après la
lumière qu'elle reçoit~; une feuille mûre ne peut pas changer.
\end{solution}

\begin{solution}{exo:b2:plant-plasticity:4}
Les \omterm{prop:b2:plant-plasticity:stomata}{cellules de garde} absorbent du potassium et de l'eau, gonflent et,
leurs parois internes étant plus épaisses, s'incurvent pour ouvrir le
pore. La lumière bleue et une faible concentration interne de
$\mathrm{CO_2}$ les ouvrent le matin~; l'\omterm{def:b2:plant-flowering:dormancy}{acide abscissique}, venu des
racines dans un sol qui s'assèche ou produit par la feuille, les ferme
en période de sécheresse, de même qu'un air très sec.
\end{solution}

\begin{solution}{exo:b2:plant-plasticity:5}
$\varphi = R/(R + 0.77\,FR)$ avec $FR = 1$~: $1.15$~: $0.60$~; $0.7$~:
$0.48$~; $0.2$~: $0.21$~; $0.05$~: $0.06$. Le seuil $0.4$ se situe entre
la trouée ($0.48$) et le couvert ($0.21$)~: l'\omterm{thm:b2:plant-plasticity:phytochrome}{évitement de l'ombre} se
déclenche sous le couvert.
\end{solution}

\begin{solution}{exo:b2:plant-plasticity:6}
$E = 0.25\times 0.025 = \qty{6.25}{mmol.m^{-2}.s^{-1}}$~; $A =
0.25\times 100/1.6 = \qty{15.6}{\micro mol.m^{-2}.s^{-1}}$~; $A/E =
2.5\times 10^{-3}$ mol de carbone par mol d'eau, c'est-à-dire
400 molécules d'eau par carbone, $400\times 18/12 = \qty{600}{g}$ d'eau
par gramme de carbone. En 12 heures~: $6.25\times 10^{-3}\times 43200 =
\qty{270}{mol/m^2}$, soit \qty{4.9}{L} par mètre carré.
\end{solution}

\begin{solution}{exo:b2:plant-plasticity:7}
$E = 0.25\times 0.005 = \qty{1.25}{mmol.m^{-2}.s^{-1}}$, $A$ inchangé~:
$A/E = 0.0125$, 80 molécules d'eau par carbone, \qty{120}{g/g} --- le
cinquième du coût diurne.
\end{solution}

\begin{solution}{exo:b2:plant-plasticity:8}
$v = 2r^{2}\Delta\rho g/9\eta = 2\times 4\times 10^{-12}\times
400\times 9.81/0.09 = \qty{3.5e-7}{m/s}$~: \qty{0.35}{\micro m/s},
environ \qty{30}{s} pour traverser la cellule. Sur un clinostat, la
direction de la gravité par rapport à la cellule change toutes les
quelques secondes de sédimentation~: les grains ne s'accumulent jamais
d'un côté, et le stimulus moyenné sur le temps d'intégration de la
plante (quelques minutes) est nul.
\end{solution}

\begin{solution}{exo:b2:plant-plasticity:9}
Extension moyenne en deux heures~: $0.04$~; faces~:
$1.2\times 0.04 = 0.048$ et $0.8\times 0.04 = 0.032$~; $\theta = 20\times 0.016/2 = 0.16$
rad, environ \qty{9}{\degree}. La courbure s'accumule à raison de $0.08$
rad par heure~; $\pi/2$ demande environ \qty{20}{h}.
\end{solution}

\begin{solution}{exo:b2:plant-plasticity:10}
Les réserves permettent \qty{50}{cm} de tige. Massif d'orties~:
\qty{30}{cm} en 7.5 jours pour \qty{30}{mg} --- la plantule atteint la
lumière avec des réserves de reste. Bois~: \qty{10}{m} sont hors de
portée~; la plantule dépense ses \qty{50}{mg} en un demi-mètre de tige
pâle en douze jours et meurt étiolée, alors qu'une stratégie de
tolérance à l'ombre --- quelques feuilles fines, une croissance lente ---
aurait pu la maintenir en vie pendant des années.
\end{solution}

\begin{solution}{exo:b2:plant-plasticity:11}
La glace se forme d'abord dans les espaces extracellulaires, où la
solution est plus diluée~; la pression de vapeur au-dessus de la glace
est plus faible qu'au-dessus de l'eau de la cellule, si bien que l'eau
quitte la cellule pour rejoindre la glace~: la cellule se déshydrate et
ses membranes s'effondrent. Une semaine de froid charge la cellule de
sucres et de protéines protectrices qui abaissent son point de
congélation et stabilisent les membranes, et fait produire des
protéines antigel qui limitent la croissance des cristaux. Maintenir le
programme actif toute l'année détourne le carbone de la croissance vers
la protection et ralentit la division~: une plante endurcie en
permanence est une petite plante.
\end{solution}

\begin{solution}{exo:b2:plant-plasticity:12}
Le signal de rouge lointain est honnête lorsque la voisine peut être
dépassée, et mensonger lorsque c'est un arbre~; la semaine douce est
honnête en avril et mensongère en février~; dans une culture dense,
toutes les voisines d'une plante sont ses égales, si bien que le signal
est honnête dans sa forme mais inutile, et les sélectionneurs retiennent
des plantes qui l'ignorent. Une \omterm{def:b2:plant-plasticity:plasticity}{norme de réaction} est la trace de la
fréquence à laquelle chaque signal a dit vrai dans le passé de la
lignée, et elle échoue chaque fois que le présent diffère de ce passé.
\end{solution}

\begin{solution}{pb:b2:plant-plasticity:1}
\textbf{1.} À découvert~: $1.15/(1.15 + 0.77) = 0.60$~; sous couvert~:
$0.2/(0.2 +
0.77) = 0.21$.
\textbf{2.} $\varphi$ faible~: des plantes au-dessus --- s'allonger. Sous
un voile d'ombrage, $\varphi$ reste à $0.60$~: c'est une journée sombre,
non une voisine~; pas d'\omterm{thm:b2:plant-plasticity:phytochrome}{évitement de l'ombre}, seulement une croissance
ralentie faute de lumière.
\textbf{3.} $1 + 2(0.60 - 0.21) = 1.78$~: \qty{2.7}{cm} par jour.
\textbf{4.} \qty{37}{cm} contre \qty{21}{cm}.
\textbf{5.} Une courbure vers la trouée (phototropisme, par le récepteur
de la lumière bleue, la phototropine) et un allongement plus rapide et
une réorientation des feuilles du côté où $R{:}FR$ est le plus élevé (par
le phytochrome).
\textbf{6.} Les \omterm{prop:b2:plant-plasticity:sunshade}{feuilles d'ombre} sont construites grandes et fines pour
capter à moindre coût une lumière rare~; en plein soleil, elles
surchauffent, subissent la photo-inhibition et brûlent, et la plante
doit les remplacer par des \omterm{prop:b2:plant-plasticity:sunshade}{feuilles de lumière}.
\textbf{7.} \qty{2.7}{cm} par jour sur une zone de \qty{15}{mm}
représentent \qty{0.11}{cm} par heure, soit une extension relative de
$0.074$ par heure~; face ombragée $1.3\times 0.074 = 0.096$, face
éclairée $0.7\times 0.074 =
0.052$.
\textbf{8.} $\theta = 15\times(0.096 - 0.052)/2 = 0.33$ rad, environ
\qty{19}{\degree} en une heure.
\textbf{9.} \qty{60}{\degree} font $1.05$ rad~: environ trois heures.
\textbf{10.} $v = 2\times 6.25\times 10^{-12}\times 400\times
9.81/0.09 = \qty{5.5e-7}{m/s}$, soit \qty{0.55}{\micro m/s}~:
\qty{22}{s} pour traverser la cellule.
\textbf{11.} Le phototropisme~; et l'\omterm{thm:b2:plant-plasticity:phytochrome}{évitement de l'ombre} accroît
lui-même l'inclinaison que tolère la pousse~: la lumière est la ressource
disputée, et une tige qui se redresserait contre la gravité retournerait
à l'ombre.
\textbf{12.} La racine se courbe vers le bas jusqu'à redevenir
verticale~: les statolithes sédimentent sur la face inférieure, l'auxine
s'y accumule et, dans une racine, l'excès d'auxine inhibe la croissance,
si bien que la face supérieure croît plus vite. Même redistribution,
sensibilité opposée des cellules.
\textbf{13.} $E = 0.2\times 0.015 = \qty{3}{mmol.m^{-2}.s^{-1}}$~; sur
$2\times 10^{-3}$ \unit{m^2} et \qty{50400}{s}~: \qty{0.30}{mol}, soit
\qty{5.4}{g} d'eau par jour.
\textbf{14.} $A = 0.2\times 80/1.6 = \qty{10}{\micro mol.m^{-2}.s^{-1}}$~;
par jour, $0.50$ mol de carbone par mètre carré~; $A/E = 3.3\times
10^{-3}$~: 300 molécules d'eau par carbone, \qty{450}{g/g}.
\textbf{15.} $0.50\times 2\times 10^{-3} = \qty{1}{mmol}$~: \qty{12}{mg}
de carbone par jour.
\textbf{16.} Trouée~: $E = \qty{6}{mmol.m^{-2}.s^{-1}}$, \qty{10.9}{g}
par jour~; efficience \qty{900}{g/g}.
\textbf{17.} $E = 0.02\times 0.03 = \qty{0.6}{mmol.m^{-2}.s^{-1}}$
(divisé par dix)~; $A = 0.02\times 200/1.6 = \qty{2.5}{\micro
mol.m^{-2}.s^{-1}}$ (divisé par quatre). L'eau diminue davantage~: à
mesure que la feuille abaisse son $\mathrm{CO_2}$ interne, le gradient
pour le carbone s'accroît, alors que celui de l'eau ne le peut pas.
\textbf{18.} \qty{4.25}{g} d'eau, dont \qty{0.85}{g} peuvent être
perdus~; à \qty{10.9}{g} par jour ($0.78$ g par heure), elle flétrit en
une heure environ.
\textbf{19.} Le métabolisme CAM est lent et exige un tissu de stockage
succulent~; une plantule qui fait la course vers la lumière sous un
couvert, où la perte d'eau est de toute façon faible, a besoin de
vitesse, non d'économie.
\textbf{20.} $100/2 = \qty{50}{cm}$.
\textbf{21.} Massif d'orties~: \qty{40}{cm} en 15 jours pour \qty{80}{mg}
--- elle atteint la lumière. Bois~: \qty{8}{m} sont inaccessibles~; elle
dépense ses réserves en un demi-mètre et meurt. Mieux valait, dans le
bois, rester courte, produire des \omterm{prop:b2:plant-plasticity:sunshade}{feuilles d'ombre} fines et attendre une
trouée --- la manière du hêtre.
\textbf{22.} Bouleau~: une norme raide, avec un fort allongement pour un
$\varphi$ faible, car les voisines d'un pionnier sont de sa taille.
Hêtre~: une norme plate, peu d'allongement et beaucoup de tolérance, car
ses voisins sont des arbres.
\textbf{23.} Désactiver la réponse d'\omterm{thm:b2:plant-plasticity:phytochrome}{évitement de l'ombre} --- maintenir
le phytochrome B actif ou supprimer les \omterm{prop:b2:cell-differentiation:transcriptional}{facteurs de transcription} qu'il
freine --- et vérifier la date de floraison, la germination et la
solidité des tiges, que contrôle la même voie.
\textbf{24.} Une plante incapable de percevoir ses voisines est dépassée
chaque fois qu'elle aurait pu gagner~; une plante qui croit toujours le
signal se gaspille sous chaque arbre. La norme gagnante parie sur le
signal en proportion de la fréquence à laquelle, dans le passé de la
lignée, il a dit vrai.
\textbf{25.} $\varphi = 0.60$ à découvert et $0.21$ sous le couvert~;
\qty{37}{cm} au bout de 14 jours~; \qty{5.4}{g} d'eau par jour à
l'ombre, \qty{10.9}{g} dans la trouée.
\end{solution}
